\documentclass[12pt]{article}
\usepackage{graphicx} % Required for inserting images
\usepackage[margin=0.5in]{geometry}
\usepackage{amsmath, amssymb}
\usepackage{mathrsfs}


\usepackage{soul}


% Dr. David Brown Commands
\newcommand{\ds}{\displaystyle}
\newcommand{\R}{\mathbb{R}}
\newcommand{\N}{\mathbb{N}}
\newcommand{\Q}{\mathbb{Q}}
\newcommand{\C}{\mathbb{C}}


% Commands to get script r
\usepackage{calligra}
\DeclareMathAlphabet{\mathcalligra}{T1}{calligra}{m}{n}
\DeclareFontShape{T1}{calligra}{m}{n}{<->s*[2.2]callig15}{}
\newcommand{\scripty}[1]{\ensuremath{\mathcalligra{#1}}}

% Unit commands
\newcommand{\degree}{\text{°}}
\newcommand{\unitSpeed}{\frac{\text{m}}{\text{s}}}
\newcommand{\unitAcceleration}{\frac{\text{m}}{\text{s}^2}}

% Constant commands
\newcommand{\avogadro}{6.022\times10^{23}}

\newcommand{\boltzmannConstK}{1.381\times10^{-23}\frac{\text{J}}{\text{K}}}
\newcommand{\boltzmannConstInEV}{8.617\times10^{-5}\frac{\text{eV}}{\text{K}}}
\newcommand{\planckConst}{6.626\times10^{-34}\text{J}\cdot\text{s}}

\newcommand{\atmP}{1.01325\times10^5\,\text{Pa}}

% Commands to easily write partial derivatives
\newcommand{\partDeriv}[2]{\frac{\partial #1}{\partial #2}}
\newcommand{\doublePartDeriv}[2]{\frac{\partial^2 #1}{\partial^2 #2}}


\title{Problem 7.9}
\author{Gabriel Decker}


\begin{document}



\maketitle
\begin{flushleft}

For this problem, we'll compute the quantum volume $v_Q$ for Nitrogen gas.
\begin{equation}
    v_Q=\left(\frac{h}{\sqrt{2\pi mkT}}\right)^3
\end{equation}
With this, we'll justifty the use of Boltzmann statistics when dealing with nonrelativistic ideal gases near room temperature by using the following condition.
The condition for when Boltzmann statistics is sufficient is when we have the following.
It comes from ensuring that the number of possible single-particle-states is much greater than the number of particles.
\begin{equation}
    \frac{V}{N}\gg v_Q
\end{equation}\\~\\

First, let's find $v_Q$.
For one molecule of Nitrogen gas, $m=4.65186\times10^{-26}\,\text{kg}$, and we're dealing with a situation where $T=300\,\text{K}$ and $P=\atmP$.
$$v_Q=\left(\frac{h}{\sqrt{2\pi mkT}}\right)^3$$
$$\implies v_Q=\left(\frac{\planckConst}{\sqrt{2\pi(4.65186\times10^{-26}\,\text{kg})(\boltzmannConstK)(300\,\text{K})}}\right)^3$$
$$\implies v_Q=6.90354\times10^{-33}\,\text{m}^3$$\\~\\

Now that we have $v_Q$, we can use it to examine the given condition.
First, recall that for an ideal gas, $V=NkT/P$.
Therefore, we get the following.
$$\frac{V}{N}\gg v_Q$$
$$\implies\frac{kT}{P}\gg v_Q$$
$$\implies\frac{(\boltzmannConstK)(300\,\text{K})}{\atmP}\gg 6.90354\times10^{-33}\,\text{m}^3$$
$$\implies4.089\times10^{-26}\,\text{m}^3\gg 6.90354\times10^{-33}\,\text{m}^3$$
This is certainly true, so Boltzmann statistics is definitely sufficient for studying this system.\\~\\

Now, let's find the required temperature (ignoring gas changing state) required to make the left comparable to the right.
We'll call that condition when the left is only ten times the right.
$$\frac{kT}{P}=10v_Q$$
$$\implies T=\frac{10\cdot Pv_Q}{k}$$
$$\implies T=\frac{10(\atmP)(6.90354\times10^{-33}\,\text{m}^3)}{\boltzmannConstK}$$
$$\implies T=0.000507\,\text{K}$$
So, quantum statistics is pretty much never needed for this system.



\end{flushleft}

\end{document}
